\documentclass[11pt,a4paper]{article}
\usepackage[utf8]{inputenc}
\usepackage[T1]{fontenc}
\usepackage{amsmath,amssymb,amsthm}
\usepackage{physics}
\usepackage{geometry}
\usepackage{hyperref}
\usepackage{graphicx}

\geometry{margin=2.5cm}

\title{Error Propagation in Measurement-Based\\
       Iterative Quantum Solvers:\\
       The Role of Post-Selection Probability}
\author{}
\date{March 2026}

\begin{document}
\maketitle

%======================================================================
\section{Introduction}
%======================================================================

This note analyses the dominant source of error in the measurement-based
(Quantum State Tomography) version of the LCU linear-pendulum solver.
We show that the finite post-selection probability \(p_{\text{succ}} < 1\)
leads to an \emph{exponential damping} of the state-vector norm when the
simulation is iterated over many time-steps, causing the quantum
trajectory to spiral inward with respect to the exact solution.

%======================================================================
\section{Setup}
%======================================================================

The system under study is the linear pendulum
\begin{equation}\label{eq:ode}
  \ddot\theta + \omega^2\,\theta = 0,
\end{equation}
discretised with the forward-Euler scheme:
\begin{equation}
  \mathbf v_{n+1}
  = A\,\mathbf v_n, \qquad
  A = I + \Delta t\begin{pmatrix} 0 & 1 \\ -\omega^2 & 0 \end{pmatrix},
  \qquad
  \mathbf v_n = \begin{pmatrix} \theta_n \\ \dot\theta_n \end{pmatrix}.
\end{equation}

In the LCU algorithm, \(A\) is decomposed as a linear combination of
unitaries
\begin{equation}
  A = c_0\,I + c_1\,(\pm X) + c_2\,(ZX),
\end{equation}
with \(c_0=1\), \(c_1=|\Delta t(1-\omega^2)/2|\),
\(c_2=\Delta t(1+\omega^2)/2\), and normalisation factor
\(\alpha = c_0+c_1+c_2\).

A 3-qubit circuit encodes the state \(\mathbf v_n\) in one system qubit
(\(q_1\)) and uses two ancilla qubits (\(q_0,q_2\)) for the LCU
preparation/un-preparation.  After post-selecting the ancillas in
\(\ket{00}\), the system qubit contains
\begin{equation}
  \ket{\psi_{\text{out}}}
  \propto A\ket{\psi_{\text{in}}}.
\end{equation}

%======================================================================
\section{Statevector (Exact) Simulation}
%======================================================================

In the \textbf{statevector} approach, the full quantum state is computed
analytically (no measurement noise).  Post-selection is carried out by
reading amplitudes directly:
\begin{equation}
  x_{n+1} = \Re\!\bigl[\alpha\,\lVert\mathbf v_n\rVert\,
             \text{amp}(\ket{000})\bigr],\qquad
  y_{n+1} = \Re\!\bigl[\alpha\,\lVert\mathbf v_n\rVert\,
             \text{amp}(\ket{010})\bigr].
\end{equation}
Since no sampling is involved, the quantum solution reproduces the
\emph{classical Euler} trajectory \textbf{exactly} (up to floating-point
errors).

%======================================================================
\section{Measurement-Based Simulation and \texorpdfstring{$p_{\text{succ}}$}{p\_succ}}
%======================================================================

When measurements with a finite number of shots \(N_s\) are used instead,
the state must be reconstructed via Quantum State Tomography (QST).
Measurements are performed in the \(Z\), \(X\), and \(Y\) bases, and
only the events where both ancillas register \(\ket{0}\) (i.e., the
post-selection condition is satisfied) are kept.

\subsection{Post-selection probability}

The probability that the ancillas are found in \(\ket{00}\) after one
LCU step is
\begin{equation}\label{eq:psucc}
  p_{\text{succ}}
  = \lVert \Pi_{00}\,U_{\text{LCU}}\ket{\psi_{\text{in}}} \rVert^2
  < 1.
\end{equation}
Physically, \(p_{\text{succ}}\) measures how ``close'' the operator
\(A/\alpha\) is to a true unitary.  For the forward-Euler step with
small \(\Delta t\), \(A\approx I\) and \(p_{\text{succ}}\) is close to
\(1\), but never exactly \(1\).

\subsection{Norm shrinkage per step}

In the exact simulation, the post-selected component has norm
\begin{equation}
  \lVert A\,\mathbf v_n\rVert / \alpha.
\end{equation}
The rescaling by \(\alpha\) \emph{exactly} compensates the LCU
normalisation, so the physical norm is preserved.

In the measurement-based version, however, the amplitudes are
\emph{estimated} from the post-selected counts.  The effective
``measured norm'' is
\begin{equation}\label{eq:measured-norm}
  \widetilde{\lVert\mathbf v_{n+1}\rVert}
  = \alpha\,\lVert\mathbf v_n\rVert\,\sqrt{\hat p_{\text{succ}}},
\end{equation}
where \(\hat p_{\text{succ}}\) is the \emph{empirical} post-selection
probability (number of successful events divided by total shots).

Because \(\hat p_{\text{succ}} < 1\), each step \textbf{contracts} the
norm by a factor \(\sqrt{\hat p_{\text{succ}}}\).

%======================================================================
\section{Cumulative Error After \texorpdfstring{$n$}{n} Steps}
%======================================================================

After \(n\) consecutive steps, the measured norm satisfies
\begin{equation}\label{eq:cum}
  \widetilde{\lVert\mathbf v_n\rVert}
  \approx
  \lVert\mathbf v_0\rVert\;
  \prod_{k=0}^{n-1}\alpha^{(k)}\sqrt{\hat p_{\text{succ}}^{(k)}}\,.
\end{equation}
If we assume that \(\alpha\) and \(\hat p_{\text{succ}}\) are roughly
constant across steps (valid for small \(\Delta t\) and slowly-varying
solutions), this simplifies to
\begin{equation}\label{eq:exp-decay}
  \widetilde{\lVert\mathbf v_n\rVert}
  \approx
  \lVert\mathbf v_0\rVert\;\bigl(\alpha\sqrt{\bar p_{\text{succ}}}\bigr)^{n},
\end{equation}
where \(\bar p_{\text{succ}}\) is the average post-selection probability.

\paragraph{Key observation.}
In the exact simulation, the factor \(\alpha\) already accounts for
the LCU normalisation and the result is
\(\lVert\mathbf v_{n+1}\rVert = \lVert A\mathbf v_n\rVert\)
(no square root of \(p_{\text{succ}}\) involved).
The culprit is the \(\sqrt{\hat p_{\text{succ}}}\) factor: even though it
is very close to \(1\), the product of \(n\) such factors decays
exponentially:
\begin{equation}
  \bigl(\sqrt{\bar p}\bigr)^{n}
  = \exp\!\Bigl(-\frac{n}{2}\ln\frac{1}{\bar p}\Bigr)
  \xrightarrow[n\to\infty]{} 0.
\end{equation}

\paragraph{Numerical example.}
With \(\Delta t=0.003\), \(\omega^2=9.8\), a typical value is
\(\bar p_{\text{succ}} \approx 0.92\).  After \(n=1000\) steps:
\[
  (\sqrt{0.92})^{1000}
  = (0.9592)^{1000}
  \approx e^{-41.6}
  \approx 10^{-18},
\]
which means the norm has essentially vanished.  In practice, the
trajectory begins to diverge visibly after only \(\sim 100\text{--}200\)
steps.

%======================================================================
\section{Additional Sources of Error}
%======================================================================

Besides norm shrinkage, additional error mechanisms include:

\begin{enumerate}
\item \textbf{Finite-sample noise in Bloch-vector components.}
      The estimators \(\hat r_z\), \(\hat r_x\), \(\hat r_y\) are
      subject to shot noise of order \(1/\sqrt{N_s}\).  These errors
      perturb the \emph{direction} of \(\mathbf v_{n+1}\), adding
      angular jitter that accumulates as a random walk.
\item \textbf{Global-sign ambiguity.}
      QST from expectation values determines the state only up to a
      global phase.  The dot-product heuristic used to resolve the sign
      may occasionally pick the wrong branch, causing a sudden
      ``phase flip'' in the trajectory.
\item \textbf{Basis-dependent post-selection rates.}
      The three bases (Z, X, Y) may have slightly different
      \(p_{\text{succ}}\) values.  Averaging them introduces a small
      systematic bias.
\end{enumerate}

%======================================================================
\section{Mitigation Strategies}
%======================================================================

\begin{itemize}
\item \textbf{Increase shot count.}  Reducing \(1/\sqrt{N_s}\) noise
      slows the angular jitter but does \emph{not} eliminate the
      \(p_{\text{succ}}\)-induced norm decay.
\item \textbf{Use smaller \(\Delta t\).}  This pushes \(A\) closer to
      \(I\), increasing \(p_{\text{succ}}\), but requires more steps for
      the same physical time.
\item \textbf{Readout error mitigation (REM).}  Calibrating and
      correcting measurement errors can improve the estimates of
      \(\hat p_{\text{succ}}\) and Bloch components.
\item \textbf{Norm renormalisation.}  Explicitly rescaling \(\mathbf v_n\)
      to its expected norm at each step partially compensates the decay,
      but requires external knowledge of the solution.
\item \textbf{Amplitude estimation.}  Replacing simple sampling with
      quantum amplitude estimation could, in principle, yield
      \(p_{\text{succ}}\) with quadratically fewer shots, but at a much
      higher circuit cost.
\end{itemize}

%======================================================================
\section{Conclusion}
%======================================================================

The divergence of the measurement-based LCU solver from the exact
(statevector) solution is \textbf{not a bug} but a fundamental
consequence of iterative post-selection with finite statistics.
The dominant effect is the exponential decay of the state-vector norm
governed by \((\sqrt{p_{\text{succ}}})^n\), which causes the quantum
trajectory to spiral inward in phase space.

For publication purposes, the statevector simulation serves as the
``ground truth'' demonstrating the correctness of the LCU algorithm,
while the measurement-based simulation illustrates the practical
challenges of running iterative quantum ODE solvers on near-term
hardware.

\end{document}
